\section{Conclusions}
\label{sec:conclusions}
A conservative finite-volume treatment of nine Pb--Bi--I gases and four
shared condensed wall inventories has been exercised in 15 baseline
and 64 follow-up cases. The reported calculations use a specialised
ideal-gas equilibrium kernel, a prescribed helium carrier and reversible
iodide wall exchange with half-cell diffusion resistance. Material
bookkeeping is strong: the worst raw source-driven follow-up element
residue is $8.85\times10^{-13}$ of supplied material. Conservation is
a numerical property of this model, rather than evidence that all
physical pathways have been represented.

The isolated time-step comparisons give iodine profile changes below
0.323\%, while medium-to-fine mesh comparisons give 2.06--10.58\%.
Stable total deposition and small centroid changes therefore do not
establish spatial convergence of local profiles. Closely competing
1 cm bins can also make a reported peak temperature switch by tens of
kelvin under refinement. Further spatial and full-run iteration checks
are required before claiming precise local predictive accuracy.

Thermodynamic and source assumptions materially affect the computed
iodide partition. On the refined steel case, a $+29$ kJ/mol BiI gas
Gibbs-energy surrogate changes the BiI$_3$ iodine share from 94.57\%
to 2.00\%. This sensitivity survives refinement but does not determine
which thermodynamic branch represents the experiment. Q1 agreement is
conditioned on inferred deposit masses, and Q2 retains a deposit-derived
iodine source and an unmeasured metal saturation.

For silica, input-derived temperature-coordinate shifts reduce the
original 4--5 cm position discrepancies to 0--1 cm. The improvement
identifies coordinate registration as a consequential uncertainty;
it does not establish a corrected experimental dataset. The supplied
paper does not resolve that registration or the flow-reference conditions.

The present evidence supports a methods and input-sensitivity study.
An independent predictive validation would require better-constrained
source and thermodynamic inputs, consistent experimental coordinates,
spatially qualified profile calculations and treatment of the
three-hour release and cooling protocol. The completed cases identify
where that additional evidence is needed without hiding the original
profile discrepancies.
